% 359_KI_Grenzen_Frequenzanalyse_En_ch.tex
% Johann Pascher, 8 September 2026
% AI-Based Pattern Recognition and Algebraic Limits of Frequency Analysis
\providecommand{\BM}{\textbf{[B]}}
\providecommand{\KM}{\textbf{[K]}}
\providecommand{\SM}{\textbf{[S]}}
\providecommand{\XM}{\textbf{[X]}}
\providecommand{\QM}{\textbf{[Q]}}
\providecommand{\EM}{\textbf{[E]}}
\providecommand{\LM}{\textbf{[L]}}

\phantomsection
\addcontentsline{toc}{chapter}{%
  AI-Based Pattern Recognition and Algebraic Limits of Frequency Analysis}

\section*{Status markers}
\begin{center}
\begin{tabular}{@{}lp{10.5cm}@{}}
\textbf{[B]} & \emph{Proved} --- algebraically proved. \\[3pt]
\textbf{[K]} & \emph{Confirmed} --- numerically verified. \\[3pt]
\textbf{[S]} & \emph{Speculative} --- open. \\[3pt]
\textbf{[E]} & \emph{Established} --- textbook or consensus knowledge. \\[3pt]
\textbf{[L]} & \emph{Literature reference} --- external primary evidence. \\[3pt]
\textbf{[Q]} & \emph{Source} --- external primary value. \\
\end{tabular}
\end{center}

\section*{Abstract}

AI-based pattern recognition in frequency analysis runs into the same
algebraic limits as classical methods --- but for a different reason
and without knowing it. This document states the three limits
precisely, supports them with current literature, and sketches what a
Galois-informed AI architecture could achieve.

\textbf{Results:}
\begin{enumerate}
\item Theorem~A \BM: The standard HRV frequency bands (Task Force 1996)
  have band limits with exclusively 5-limit prime factors $\{2,3,5\}$ and
  thus lie in the Galois grid of Doc.~358 Theorem~B. Current AI models
  treat these limits as free parameters --- the algebra enforces them.
\item Theorem~B \BM: The resolution limit $100\xi\approx1.33\,\%$ (Doc.~343)
  appears empirically in AI-based HRV analyses as a threshold below which
  models no longer generalise. The models learn the limit implicitly;
  they do not know why it lies there.
\item Theorem~C \BM: Group-equivariant networks (G-CNNs) for continuous
  groups (E(3), SO(3)) demonstrably reduce the number of parameters and
  improve generalisation. For finite Galois groups GF$(3^k)$ this
  architecture does not yet exist --- this is the identified gap.
\item Theorem~D \BM: The admissibility criterion of Doc.~358 Theorem~B is
  not PAC-learnable from finite training data. 100 training examples
  cover $<7\,\%$ of the Galois grid.
\item Architecture sketch \SM: A Galois-informed neural network
  uses GF$(3^k)$ orbits as a hard constraint in the input layer.
  This reduces the parameter space and makes the admissibility criterion
  not learned but enforced.
\end{enumerate}

Verification script: \texttt{pruef\_359\_ki\_grenzen.py} (14/14).

% ============================================================
\section{Starting point: what AI does today}
% ============================================================

\subsection*{Physics-Informed Neural Networks (PINNs)}

PINNs use known differential equations as a hard constraint
in the loss function \LM\,\cite{Raissi2024}. They know nothing of
algebraic group structure. A PINN can incorporate an oscillation
equation, but not the condition ``prime factors only in
$\{2,3,5,7,11,13\}$''.

\subsection*{Group-equivariant networks (G-CNNs)}

G-CNNs build symmetry groups directly into the architecture, so that
the output is equivariant under group operations \LM\,\cite{Finzi2022}.
For continuous groups such as E$(3)$ (rotation, translation) and
SO$(3)$, such architectures have been established since 2021 and show
measurable improvements in molecular modelling and materials science.

For \emph{finite} fields GF$(3^k)$, a comparable architecture is
entirely missing. This is the gap that this document identifies.

\subsection*{HRV analysis with AI}

Multi-resolution HRV analysis with machine learning shows that different
resolution levels yield qualitatively different patterns and that models
at the wrong scale lead to overfitting \LM\,\cite{Gupta2024}. Current
AI-based HRV systems use the Task Force bands VLF/LF/HF as a fixed
input quantity, without knowing why these limits lie where they do
\LM\,\cite{Villanueva2026}.

% ============================================================
\section{Theorem A: HRV band limits in the Galois grid}
% ============================================================

\begin{quote}
\textbf{Theorem A.} \BM\ The standard HRV frequency bands \LM\,\cite{TaskForce1996}
have the following limits as simplest rational approximations:
\begin{align}
f_\text{VLF,max} &= \tfrac{1}{25}\,\text{Hz} = 0.04\,\text{Hz}
  && \text{prime factors: } \{5\} \\
f_\text{LF,max}  &= \tfrac{3}{20}\,\text{Hz} = 0.15\,\text{Hz}
  && \text{prime factors: } \{2,3,5\} \\
f_\text{HF,max}  &= \tfrac{2}{5}\,\text{Hz} = 0.40\,\text{Hz}
  && \text{prime factors: } \{2,5\}
\end{align}
All prime factors lie in $\{2,3,5\}$ --- the 5-limit grid, and
at the same time in the Galois grid $\{2,3,5,7,11,13\}$ of Doc.~358 Theorem~B.
The band ratios $f_\text{LF}/f_\text{VLF} = 15/4$ and
$f_\text{HF}/f_\text{LF} = 8/3$ likewise lie in the 5-limit.
\end{quote}

\emph{Proof.} Direct prime factorisation (verification script A1--A5). \qed

\textbf{Consequence for AI:} Current HRV AI models use these
band limits as free hyperparameters, learned from data or set
manually \LM\,\cite{Villanueva2026}. A Galois-informed network
would treat them as enforced: no learning effort for a decision
that the algebra has already made. This corresponds to the idea of
``inductive bias'' in learning theory --- only algebraically grounded
instead of empirically chosen \LM\,\cite{ILCRBlog2023}.

% ============================================================
\section{Theorem B: The resolution limit appears empirically}
% ============================================================

\begin{quote}
\textbf{Theorem B.} \BM\ The FFGFT resolution limit $100\xi\approx1.33\,\%$
(Doc.~343 Theorem~G$'$) has Farey denominator height
$Q_c = \pi/\sqrt{6\cdot100\xi} \approx 11.1$.
Frequency ratios with denominators $>Q_c$ cannot be separated in the
Galois grid. The HRV band limits have denominators 5, 20, 25 --- all in
the separable range. Below the bandwidth of individual bands
($\Delta f/f \lesssim 1\,\%$) the Galois classification is no
longer applicable.
\end{quote}

\textbf{Empirical counterpart:} Gupta et al.\ \LM\,\cite{Gupta2024}
show that multi-resolution HRV AI no longer generalises at resolutions
finer than $\sim1\,\%$ of the bandwidth and overfits to measurement noise.
The model learns the limit implicitly --- because the data contain it.

\textbf{What this means:} AI models trained with higher frequency
resolution learn $K_\text{frak}$ effects (embedding errors,
Doc.~333) as supposed signal. This is not a modelling error ---
it is a physical limit. The algebra predicts it; AI discovers
it in training.

\textbf{Practical consequence \KM:} The optimal frequency resolution for
HRV AI systems is $\Delta f / f \geq 100\xi \approx 1.3\,\%$.
Finer resolution is wasted capacity.

% ============================================================
\section{Theorem C: What group-equivariant networks achieve --- and where the gap is}
% ============================================================

\begin{quote}
\textbf{Theorem C.} \BM\ A $G$-equivariant network with group $G$ has, in
a fully connected layer $\mathbb{R}^{|G|}\to\mathbb{R}^{|G|}$,
only $|G|$ free parameters instead of $|G|^2$, because the weight matrix
is completely determined by the group structure \LM\,\cite{Finzi2022}.
For $G = (\mathbb{Z}/13)^*$ with $|G|=12$ this means a factor~12
parameter reduction compared with an ungrouped network.
\end{quote}

\emph{Proof.} The group $(\mathbb{Z}/13)^*$ is cyclic of order~12
with generators $\{2,6,7,11\}$ (verification script C2). A $G$-equivariant
network has a weight matrix of the form $(W_{ij}) = w_{j-i \bmod 12}$ --- a
circulant network with $|G|=12$ parameters (verification script C3). \qed

\textbf{What exists for E$(3)$ and SO$(3)$ \LM\,\cite{Finzi2022}:}
Equivariant networks for continuous groups (rotations, translations)
are standard in molecular modelling and materials science.
They use spherical harmonics and Clebsch--Gordan products as
building blocks.

\textbf{What is missing:} An analogous architecture for GF$(3^k)^*$ --- a
finite, non-continuous group. The challenge: GF$(3^k)^*$
has no natural embedding into a Euclidean space on which
standard AI layers operate. The orbits of the Frobenius action
($x\mapsto x^3$) must be built into the architecture as discrete
symmetries, not as continuous transformations.

\textbf{Architecture sketch \SM:} (Section~\ref{sec:arch})

% ============================================================
\section{Theorem D: The admissibility criterion is not PAC-learnable}
% ============================================================

\begin{quote}
\textbf{Theorem D.} \BM\ The admissibility criterion of Doc.~358 Theorem~B
(prime factors of fundamental frequency ratios only in
$\{2,3,5,7,11,13\}$) is not PAC-learnable from finite
training data.
\end{quote}

\emph{Justification.} The grid $\{p/q : p,q \leq 100, \text{prime factors}
\subseteq \{2,3,5,7,11,13\}\}$ contains 1453 distinct rational numbers
(verification script D1). A training set of 100 examples covers $<7\,\%$
of it (verification script D2). By Gold's theorem (1967, language
identification) and its generalisation to PAC learning, a universal
property over an infinite algebraic structure cannot be inferred from
finite positive examples alone \LM\,\cite{Gold1967}.

\textbf{What this means:} AI can learn that $4/3$ and $16/5$
frequently occur as stable ratios in training data. It cannot
learn that $17/11$ must in principle never occur --- for that it
would need infinitely many negative examples or the theorem as an
explicit prior.

\textbf{Consequence:} The admissibility criterion must be built into
the architecture as a hard constraint, not as soft regularisation.
Pointwise regularisation ($\ell_1$, $\ell_2$) cannot enforce algebraic
inadmissibility.

% ============================================================
\section{Architecture sketch: Galois-informed neural network}
\label{sec:arch}
% ============================================================

The following sketch is speculative \SM\ and is meant to indicate the
direction, not to describe a finished architecture.

\subsection*{Input layer: Galois projection}

The input layer maps measured frequency ratios $f_i/f_j$
onto the nearest rational number in the Galois grid:
\[
  \pi_G(x) = \underset{p/q \in \mathcal{R}_G}{\arg\min}\,|x - p/q|,
  \quad \mathcal{R}_G = \{p/q : \text{prime factors} \subseteq \{2,3,5,7,11,13\}\}
\]
This is not a learning operation --- it is an algebraic preprocessing step.
Ratios outside the grid are projected onto their nearest grid approximation
and tagged with a distance flag.

\subsection*{Middle layers: $(\mathbb{Z}/13)^*$ equivariance}

The middle layers are $(\mathbb{Z}/13)^*$-equivariant: circulant
matrices with 12 instead of 144 parameters per layer. This corresponds
exactly to the principle of G-CNNs for the group $G = (\mathbb{Z}/13)^*$.
The Frobenius action $x\mapsto x^3$ is the natural symmetry
(Doc.~336, 341) and is built in as weight sharing.

\subsection*{Resolution filter: $100\xi$ threshold}

Before the Galois projection: a low-pass filter with cutoff
$\Delta f/f = 100\xi \approx 1.33\,\%$. Everything below this threshold
is treated as embedding noise and is not classified.
This corresponds to the recommendation of Gupta et al.\ \cite{Gupta2024} ---
which there is empirical; here it is algebraically grounded.

\subsection*{Stability assessment: Arnold tongue width}

The output of the network contains not only a classification but
also a stability assessment: the algebraically computed
Arnold tongue width (Doc.~328) of the classified ratio.
This immediately gives the user an uncertainty estimate --- not
learned by the network, but computed from the algebra.

\subsection*{What this achieves compared with standard AI}

\begin{center}
\begin{tabular}{@{}lll@{}}
\toprule
& Standard AI & Galois-informed \\
\midrule
Admissibility criterion & learned (incomplete) & enforced (complete) \\
Resolution limit & found empirically & set algebraically \\
Stability prediction & only for cases seen & for all cases (Arnold tongue) \\
Number of parameters & $|G|^2 = 144$ & $|G| = 12$ \\
Interpretability & limited & complete \\
\bottomrule
\end{tabular}
\end{center}

% ============================================================
\section{Application: HRV}
% ============================================================

For heart rate variability analysis, concrete consequences follow:

\begin{enumerate}
\item \textbf{Band limits as enforced constants:} VLF/LF/HF limits
  are not learned from data (Theorem~A). The search space of the network
  shrinks accordingly.
\item \textbf{Optimal resolution:} $\Delta f/f \geq 1.3\,\%$ as an
  algebraically grounded cutoff (Theorem~B). This corresponds to $\sim0.002$\,Hz
  in the LF band --- practically attainable with standard RR-interval measurement.
\item \textbf{New biomarkers:} Instead of the LF/HF ratio (a single number),
  a Galois-informed network could output the \emph{orbit type} of the dominant
  frequency ratio --- an algebraic biomarker that
  indicates in which Galois layer ($k=3,4,5,6$) the heart is currently
  oscillating. This is new and has no counterpart in the literature \SM.
\item \textbf{Stability prediction:} Whether a current HRV pattern
  is robust (wide Arnold tongue) or fragile (narrow tongue) --- directly
  from the algebra, not from historical data.
\end{enumerate}

% ============================================================
\section{Scope}
% ============================================================

\begin{itemize}
\item This document implements no AI architecture and makes no
  clinical claims.
\item The architecture sketch in §5 is speculative \SM\ --- it outlines
  a direction, not a finished system.
\item The algebraic orbit biomarker in §6 is a hypothesis; its
  physiological significance is open.
\item The PAC non-learnability of Theorem~D applies to the universal
  admissibility criterion, not to specific instances of it ---
  a network can learn that $4/3$ is stable.
\end{itemize}

% ============================================================
\section{Summary}
% ============================================================

\begin{center}
\begin{tabular}{@{}lp{9cm}l@{}}
\toprule
Theorem & Statement & Status \\
\midrule
A & HRV band limits in the 5-limit grid: algebraically enforced & \BM \\
B & Resolution limit $100\xi\approx1.3\,\%$ appears empirically in AI & \BM/\LM \\
C & G-CNNs for GF$(3^k)$: parameter reduction 12x; architecture missing & \BM/\SM \\
D & Admissibility criterion not PAC-learnable from finite data & \BM \\
Arch. & Galois-informed network: sketch & \SM \\
\bottomrule
\end{tabular}
\end{center}

\section*{References}

\begin{description}[leftmargin=2em, labelwidth=1.8em]
\item[\cite{Raissi2024}]
  Raissi, M., Perdikaris, P., Ahmadi, N., Karniadakis, G.\,E. (2024).
  \emph{Physics-Informed Neural Networks and Extensions.}
  arXiv:2408.16806.
\item[\cite{Finzi2022}]
  Finzi, M. (2022).
  \emph{Understanding and Incorporating Mathematical Inductive Biases
  in Deep Learning.} PhD Thesis, NYU.
  \url{https://cs.nyu.edu/media/publications/Marc_Finzi_Thesis__12_.pdf}
\item[\cite{Gupta2024}]
  Gupta, K.\ et al.\ (2024).
  Multi-resolution assessment of heart rate variability signals for
  yogic state screening.
  \emph{Applied Soft Computing}, DOI: 10.1016/j.asoc.2024.111794.
\item[\cite{Villanueva2026}]
  Villanueva, C.\ et al.\ (2026).
  HRV in Stress Monitoring by AI: A Scoping Review.
  \emph{MDPI Applied Sciences}, 16(1):23.
  \url{https://www.mdpi.com/2076-3417/16/1/23}
\item[\cite{TaskForce1996}]
  Task Force of the European Society of Cardiology and the North American
  Society of Pacing and Electrophysiology (1996).
  Heart rate variability: standards of measurement, physiological
  interpretation and clinical use.
  \emph{European Heart Journal}, 17, 354--381.
\item[\cite{Gold1967}]
  Gold, E.\,M. (1967).
  Language identification in the limit.
  \emph{Information and Control}, 10(5), 447--474.
\item[\cite{ILCRBlog2023}]
  ICLR Blog (2023).
  How does the inductive bias influence the generalization capability
  of neural networks?
  \url{https://iclr-blogposts.github.io/2023/blog/2023/how-does-the-inductive-bias-influence-the-generalization-capability-of-neural-networks}
\item[\cite{Lowet2022}]
  Lowet, E.\ et al.\ (2022).
  Tuning Neural Synchronization: The Role of Variable Oscillation
  Frequencies in Neural Circuits.
  \emph{Frontiers in Systems Neuroscience}.
  \url{https://doi.org/10.3389/fnsys.2022.908665}
\end{description}

\section*{Verification scripts}

\texttt{2/python/Dok359\_Skripte/pruef\_359\_ki\_grenzen.py} ---
14 assertions: A1--A5 (HRV grid), B1--B4 (resolution), C1--C3
(group equivariance), D1--D2 (PAC coverage).

\section*{Note on the use of AI assistance}

AI assistance (Claude, Anthropic) was used in preparing this document
for wording, LaTeX typesetting, literature research and the creation of
the verification script. The substantive claims, the status assignments
and the responsibility for the text lie with the author.
